% Section 06: Results and Analysis
% Structure: sRQ-organized with Delphi round progression within each sRQ

\section{Results and Analysis}\label{sec:results}

\subsection*{Introduction}

This section presents findings from the three-round Delphi study
(Section~\ref{sec:methodology}), organized by sub-research question. Within
each subsection, findings follow the round progression to show how practitioner
consensus evolved. The following subsections detail the results:

\localtableofcontents\clearpage

%% ================================================================
%% 6.2 Panel Overview
%% ================================================================
\subsection{Delphi Panel and Convergence Overview}\label{subsec:panel-overview}

This subsection describes the expert panel composition across the three Delphi
rounds and summarizes each round's procedural outcomes.

Table~\ref{tab:panel-participation} presents the panel participation summary.
Eight Enterprise Architecture (EA) practitioners completed asynchronous discovery workbooks in Round~1
(January 2026). Ten of twelve invited participants joined the synchronous
Round~2 session (February 2026). Seven participants returned for the final
synchronous Round~3 session (March 2026). For Round~2, the panel was expanded to twelve invitees (five returning, seven
new) to introduce perspectives not anchored to Round~1 submissions. For
Round~3, all seven participants were returning from earlier rounds, ensuring
continuity in the final convergence phase.

\begin{table}[!ht]
    \centering
    \caption{Panel Participation Summary}\label{tab:panel-participation}
    \begin{tabular}{|l|c|c|c|}
        \hline
        & \textbf{Round~1} & \textbf{Round~2} & \textbf{Round~3} \\
        \hline
        Format & Async workbook & Sync 120~min & Sync 120~min \\
        \hline
        Date & January 2026 & 4 February 2026 & 4 March 2026 \\
        \hline
        Invited & 8 & 12 & 10 \\
        \hline
        Participated & 8 & 10 & 7 \\
        \hline
        Platform & Word workbook & Zoom\footnote{\url{https://zoom.us/}} + MeetingWizard\footnote{\url{https://www.meetingwizard.nl/en/}} & Zoom + MeetingWizard \\
        \hline
        Voted & -- & 10 & 7 \\
        \hline
    \end{tabular}
\end{table}

Participants held diverse architectural roles including enterprise architect,
solution architect, platform engineer, enterprise data architect, software
architect, and freelance architectural consultant. This role diversity reflects
how Reference Architecture (RA) development unfolds in practice across
organizational boundaries, consistent with the broad ``Enterprise Architecture (EA) practitioners''
framing established in Section~\ref{subsec:expert-panel}.

\textbf{Round~1} focused on open discovery. Eight participants documented
19 distinct Generative AI (GenAI) practices, nine critical
incidents, and twelve challenges with workarounds. No researcher-imposed
categories constrained what participants could report, consistent with the
exploratory Delphi orientation described in Section~\ref{sec:methodology}.

\textbf{Round~2} shifted to pattern validation. The session presented six
practice categories synthesized from Round~1 analysis. Participants voted on
category importance, discussed boundary conditions, and proposed refinements. A
seventh category (Governance and Compliance Checking) emerged during live
discussion and was validated in real time. The session produced seven validated
practice categories with importance ratings and convergence data on cross-cutting
patterns.

\textbf{Round~3} targeted remaining evidence gaps. The synchronous session
executed Steps 1--14 of the planned 16-step protocol. Steps 1--12 proceeded as
designed. Steps 13--14 (artifact discussion and guidance element voting) were
compressed due to time constraints. Step~15 (Future Perspectives brainstorm) was
not executed. Step~16 (SP2 Maturity Progression vote) was collected
asynchronously after the session, with all seven participants responding.
\textit{Data Limitations} (Section~\ref{subsec:data-limitations}) addresses the implications of these
procedural variations.

%% ================================================================
%% 6.3 sRQ1: GenAI Practices in RA Development
%% ================================================================
\subsection{sRQ1: GenAI Practices in RA Development}\label{subsec:srq1}

This subsection addresses sRQ1: \textit{``What GenAI practices are EA
Practitioners currently using when developing RAs?''} We present practice
discovery from Round~1, category validation from Round~2, and cross-cutting
patterns that emerged across both rounds. These findings form a two-level
taxonomy: seven practice categories that describe task domains, and five
cross-cutting patterns that describe behavioral universals transcending any
single category.

\subsubsection{Practice Discovery (Round~1)}\label{subsubsec:practice-discovery}

Round~1 yielded 19 distinct GenAI practices from eight participants, each
documenting two to three practices with step-by-step descriptions, genesis
stories, and effectiveness assessments. The Round~1 coded analysis applied four coding dimensions: interaction patterns,
distribution of cognitive tasks, engagement quality, and practice reconfiguration.
These dimensions emerged through the open and axial coding process described in
Section~\ref{sec:methodology}, following established qualitative coding procedures
\parencite{charmaz_constructing_2014, patton_qualitative_2015}.

The practices concentrated in three RA lifecycle phases. Phase~4 (Target State
Design) attracted the most practices, with seven of 19 mapped primarily to
Research \& Exploration activities. Phase~6 (Documentation \& Communication)
accounted for six practices across First Draft \& Documentation and Stakeholder
Communication. Phase~2 (Current State Analysis) contributed two practices in
System/Codebase Analysis. Phases~1, 3, 5, 7, and~8 received no primary practice
mappings, though some practices spanned multiple phases.

During the Round~2 session, the researcher clarified the breadth of RA work to
participants. The researcher framed this as: ``it can be a pattern, it can be a
module\ldots{} and you are not necessarily just creating, sometimes you are
absorbing, going through it, reading, getting context'' (R2 Transcript,
researcher facilitation). The aim was to capture practitioner activities beyond artifact origination.

Interaction patterns varied across the panel. Research Assistant mode dominated,
with seven of 19 practices involving GenAI as an information-gathering tool.
First Draft Generator mode appeared in six practices where GenAI produced initial
content for substantial human editing. Two participants operated in Cyborg mode
with deep tooling integration via MCP servers, hooks, and skills. Two practices
reflected a Junior Colleague framing with conversational, supervisory
interaction. No participant delegated architectural decisions entirely to GenAI.

A key finding from Round~1 concerned the distribution of cognitive tasks. All eight
participants delegated information gathering to GenAI while retaining judgment
under ambiguity. Seven of eight described a hybrid synthesis pattern where GenAI
drafts content and the human reconciles, validates, and edits. Organizational
context (team dynamics, political constraints, historical decisions) remained
exclusively human-retained across the entire panel.

\subsubsection{Practice Category Validation (Round~2)}\label{subsubsec:practice-validation}

Round~2 validated the practice categories through collective deliberation. The
researcher presented six categories synthesized from Round~1 analysis. During the
flipboard discussion, participants identified a gap: governance and compliance
checking was not represented in the original six categories. The panel proposed
and validated a seventh category in real time, demonstrating the emergent nature
of the Delphi process.

Table~\ref{tab:practice-importance} presents the importance ratings for all seven
validated categories. System/Codebase Analysis received the highest mean rating
with the lowest variability, indicating strong consensus on its importance for
RA work. By contrast, Stakeholder Communication and Governance/Compliance showed
the highest variability, reflecting context-dependent value across
organizational settings.

\begin{table}[!ht]
    \centering
    \caption{Practice Category Importance Ratings (Round~2, $N = 10$)}\label{tab:practice-importance}
    \begin{tabular}{|p{6.5cm}|c|c|c|c|}
        \hline
        \textbf{Practice Category} & \textbf{$M$} & \textbf{$SD$} & \textbf{Var.\%} & \textbf{Abst.} \\
        \hline
        System / Codebase Analysis & 4.00 & 0.82 & 38 & 0 \\
        \hline
        Research \& Exploration & 3.60 & 1.17 & 55 & 0 \\
        \hline
        Adversarial Review / Stress Testing & 3.60 & 0.97 & 45 & 0 \\
        \hline
        Governance / Compliance Checking & 3.50 & 1.27 & 60 & 0 \\
        \hline
        First Draft \& Documentation & 3.10 & 1.10 & 52 & 0 \\
        \hline
        Stakeholder Communication & 3.10 & 1.37 & 65 & 0 \\
        \hline
        Deep Workflow Integration & 2.89 & 1.17 & 54 & 1 \\
        \hline
    \end{tabular}
\end{table}

The consensus--variability pattern across categories showed a clear structure.
System/Codebase Analysis exhibited both the highest importance and lowest
variability, so directive guidance is feasible for this practice. Categories with higher
variability (Stakeholder Communication, Governance/Compliance) call for
guidance that remains contextual and adaptive. Deep Workflow Integration received
the lowest mean rating and one abstention, consistent with its dependence on
individual maturity and organizational tooling policies.

\subsubsection{Cross-Cutting Patterns}\label{subsubsec:cross-cutting}

The sRQ1 findings comprise a two-level taxonomy. The seven practice categories
presented above constitute the first level: they describe \emph{what}
practitioners do, organized by task domain. Each category is
context-dependent; its importance varies with organizational setting, as the
consensus--variability profiles in Table~\ref{tab:practice-importance} confirm.
The five cross-cutting patterns below constitute the second level: they
describe \emph{how} practitioners approach GenAI work, regardless of task
domain. The Round~1 coded analysis identified these patterns through constant
comparison across participants \parencite{charmaz_constructing_2014}; Round~2
discussions reinforced them.

\textbf{Universal Verification Mandate (8/8 participants).} All eight Round~1
participants described explicit verification practices. No participant accepted
GenAI outputs without review. Verification strategies ranged from multi-round
review cycles to line-by-line comparison with source documents, from
knowledge-based assessment to treating all outputs as hypotheses. As P4
described: ``Everything is treated as hypothesis, never conclusion'' (R1
Workbook). P2 reported verifying ``60--70\% of specific claims'' against
primary sources (R1 Workbook). This universal pattern mapped to what we term \emph{engaged
augmentation}: a coding dimension that emerged from the Round~1 analysis to
distinguish active oversight from passive delegation. Engaged augmentation
characterized the entire panel.

\textbf{Compression Not Generation (5/8 participants).} Five participants framed
GenAI's value as compressing existing information rather than creating new
architectural knowledge. As P2 described: ``The value is not in generation
but in compression. It compresses reading time for surveys'' (R1 Workbook).
P3 reinforced this framing: ``The thinking is still mine; the AI accelerates
expression'' (R1 Workbook). This pattern reframes the distribution of cognitive tasks
\parencite{hutchins_cognition_1995}; GenAI handles information gathering while
humans retain sensemaking.

\textbf{Frontier Learning Through Failure (6/8 participants).} Six participants
described learning GenAI's limitations through errors. P1 caught a retention
period change from seven to five years only through verification. P2
identified a hallucinated benchmark that nearly entered a planning document.
P3 encountered bounded context suggestions that were ``technically logical but
organisationally nonsensical'' (R1 Workbook). These incidents demonstrated how
practitioners discovered the boundary between tasks inside and outside GenAI
capability through direct experience.

\textbf{Organizational Context as Hard Boundary (7/8 participants).} Seven
participants identified organizational and sociotechnical factors as outside
GenAI capability. P3 observed that ``organisational context is nearly
impossible to convey: team dynamics, political constraints, historical baggage''
(R1 Workbook). P4 noted that GenAI ``is not valuable for the sociotechnical
aspects that determine whether an architecture will actually succeed'' (R1
Workbook). This pattern identified a category of knowledge that practitioners
consistently retained rather than delegated.

\textbf{Role Identity Preservation (8/8 participants).} All eight participants
viewed GenAI as an efficiency tool, not a replacement for professional judgment.
P4 stated: ``I did this work for years without GenAI. Efficiency would
decrease but effectiveness would not'' (R1 Workbook). P3 described that he
``would be slower but functional. The thinking is still mine'' (R1 Workbook).
This pattern indicated that professional identity remained intact despite
significant workflow changes. These cross-cutting patterns, particularly the
verification mandate and role identity preservation, set the stage for
understanding how EA practice itself changed with GenAI adoption.

%% ================================================================
%% 6.4 sRQ2: Changes to EA Practice
%% ================================================================
\subsection{sRQ2: Changes to EA Practice}\label{subsec:srq2}

sRQ2 asks: \textit{``What changes to EA practice are EA practitioners
experiencing when RAs are developed using GenAI?''} The findings below cover
practice changes from Rounds~1--2, cross-role dynamics that emerged across all
three rounds, and a convergence assessment.

\subsubsection{Practice Changes (Rounds~1--2)}\label{subsubsec:practice-changes}

Round~1 identified four dimensions of practice change:

\begin{enumerate}
    \item \textbf{Creator to curator and judge.} Six of eight participants
    reported that their primary contribution moved from producing content to
    evaluating, editing, and validating AI-generated content.
    \item \textbf{New skills emerged.} Prompt engineering, output evaluation, and
    trust calibration appeared across the panel as competencies that did not exist
    before GenAI adoption.
    \item \textbf{Certain skills used less frequently.} Cold-start writing,
    exhaustive manual search, and cover-to-cover documentation reading declined.
    \item \textbf{Artifact formats adapted.} Several practitioners shifted toward
    machine-readable formats such as Markdown and Mermaid diagrams.
\end{enumerate}

Round~2 quantified these changes through significance voting.
Table~\ref{tab:change-significance} presents the results. Verification burden
shift received the highest rating with the lowest variability, indicating strong
consensus that practitioners spend less time creating and more time validating.
Compression not generation received the most contested rating, suggesting that
not all practitioners frame GenAI's contribution as purely compressive.

\begin{table}[!ht]
    \centering
    \caption{Change Significance Ratings (Round~2, $N = 10$)}\label{tab:change-significance}
    \begin{tabular}{|p{7cm}|c|c|c|}
        \hline
        \textbf{Change Theme} & \textbf{$M$} & \textbf{$SD$} & \textbf{Var.\%} \\
        \hline
        Verification burden shift & 4.40 & 0.70 & 33 \\
        \hline
        Skills used less & 3.90 & 0.57 & 26 \\
        \hline
        Role evolution (creator $\rightarrow$ curator) & 3.80 & 1.14 & 53 \\
        \hline
        Compression not generation & 2.90 & 0.99 & 47 \\
        \hline
        New skills emerging & 2.90 & 1.20 & 56 \\
        \hline
    \end{tabular}
\end{table}

The low rating for ``new skills emerging'' ($M = 2.90$) alongside the high
rating for ``verification burden shift'' ($M = 4.40$) revealed an interesting
tension. Practitioners acknowledged that their work had changed substantially
(verification) but were less inclined to characterize this change as requiring
fundamentally new skills. This pattern indicated that practitioners treated verification as an
extension of existing professional expertise, not a novel competency.

\subsubsection{Cross-Role Dynamics (Rounds~1--2--3)}\label{subsubsec:cross-role}

Cross-role dynamics emerged as an unanticipated theme across all three rounds.
In Round~1, several practitioners described GenAI-enabled practices crossing
traditional role boundaries: developers producing architectural artifacts
bottom-up, solution architects generating decision records from meeting
transcripts, and platform engineers contributing to enterprise-level governance.
Round~2 reinforced these observations. The panel's role diversity surfaced
inside-out workflows where developer-level GenAI use shapes EA artifacts before
formal governance awareness.

Round~3 probed these dynamics through the ``Ghost Architect'' provocation:
``If anyone with GenAI can produce architectural artifacts, who is the
architect?'' Seven participants contributed sixteen brainstorm responses across
three probes (bottom-up workflows, traceability, governance friction). One
participant observed:
``The friction is real, and honestly, it's almost hard to admit. You spend 15
years building up knowledge\ldots{} And then someone sits down with GenAI and
produces something that looks pretty much the same in under an hour'' (R3
Brainstorm). Another noted that ``everybody can generate an architecture now,
which is cool. But the one taking responsibility of the decision is in the end
the architect'' (R3 Brainstorm). These responses captured a shared tension:
GenAI democratized artifact production, but accountability for architectural
decisions remained with the architect.

Table~\ref{tab:cross-role} presents the cross-role pattern significance ratings
from Round~3. Collective decision capture (meeting transcripts and requirements
transformed into architecture decision records) received the highest mean rating
($M = 3.86$, $SD = 1.68$). However, all six patterns exhibited high variability
(37--79\%), indicating substantial disagreement about the significance of
cross-role shifts. Governance patterns (inside-out and outside-in) received the
lowest ratings, suggesting that GenAI had not yet materially affected governance
structures in most participants' organizations.

\begin{table}[!ht]
    \centering
    \caption{Cross-Role Pattern Significance (Round~3, $N = 5$--$7$)}\label{tab:cross-role}
    \begin{tabular}{|p{6.5cm}|c|c|c|c|}
        \hline
        \textbf{Pattern} & \textbf{$N$} & \textbf{$M$} & \textbf{$SD$} & \textbf{Var.\%} \\
        \hline
        Collective decision capture & 7 & 3.86 & 1.68 & 77 \\
        \hline
        Top-down code generation & 7 & 3.43 & 1.72 & 79 \\
        \hline
        Role boundary blurring & 7 & 2.86 & 1.57 & 72 \\
        \hline
        Bottom-up artifact generation & 7 & 2.57 & 1.27 & 58 \\
        \hline
        Outside-in governance\textsuperscript{*} & 5 & 2.40 & 1.14 & 50 \\
        \hline
        Inside-out governance & 7 & 2.00 & 0.82 & 37 \\
        \hline
        \multicolumn{5}{l}{\small \textsuperscript{*}Five of seven participants rated this item; two abstained, likely because this} \\
        \multicolumn{5}{l}{\small emergent pattern (proposed mid-session) was not applicable to their organizational context.} \\
        \hline
    \end{tabular}
\end{table}

The high variability across all cross-role patterns reflected divergent
organizational realities. Participants from organizations with mature
governance structures reported minimal boundary shifts; those in more fluid
environments described significant role redistribution. Cross-role dynamics
appeared to depend on organizational factors rather than on GenAI adoption
itself.

\subsubsection{Convergence Assessment}\label{subsubsec:convergence}

Across the three rounds, several findings stabilized while others remained
contested, consistent with the iterative convergence expected in Delphi designs
\parencite{von_der_gracht_consensus_2012}. The seven practice categories validated in Round~2 showed no
challenges or revisions in Round~3, indicating stable consensus on the practice
taxonomy. The universal verification mandate (8/8 in Round~1) was reinforced
through Round~3's heuristic refinement, where cross-validation received the
highest importance rating. The verification burden shift achieved the strongest
consensus of any change theme ($M = 4.40$, variability 33\%).

The SP2 maturity progression question, which could not be addressed during the
synchronous session, was resolved through asynchronous data collection. The panel
strongly agreed that practitioners naturally progress from basic to sophisticated
GenAI integration ($M = 4.71$, variability 22\%). At the same time, context
received moderate agreement as a factor ($M = 3.86$, variability 49\%), while the
proposition that short experience can equal long experience was rejected
($M = 2.86$, variability 31\%). The panel moderately agreed that maturity should
be measured per practice category rather than per individual ($M = 3.86$,
variability 49\%). Table~\ref{tab:maturity-progression} presents the full
results. Taken together, these ratings support a ``both/and'' interpretation:
practitioners do progress over time, but progression is practice-specific and
pace-modulated by organizational context.

\begin{table}[!ht]
    \centering
    \caption{SP2 Maturity Progression Ratings (Async, $N = 7$)}\label{tab:maturity-progression}
    \begin{tabular}{|p{7.5cm}|c|c|c|}
        \hline
        \textbf{Statement} & \textbf{$M$} & \textbf{$SD$} & \textbf{Var.\%} \\
        \hline
        Practitioners naturally progress from basic GenAI use to more sophisticated integration over time & 4.71 & 0.45 & 22 \\
        \hline
        Sophistication depends more on context than on individual experience & 3.86 & 0.99 & 49 \\
        \hline
        A practitioner with 6 months of experience can be as effective as one with 2 years, given the right context and guidance & 2.86 & 0.64 & 31 \\
        \hline
        Maturity should be measured per practice category, not per individual & 3.86 & 0.99 & 49 \\
        \hline
    \end{tabular}
\end{table}

Two areas remained empirically unresolved. First, the cross-role patterns showed
variability too high to establish consensus, suggesting these dynamics are still
emerging. Second, the relative importance of new skills versus existing expertise
remained contested, with practitioners split on whether GenAI required
fundamentally new competencies or merely extended existing professional judgment.
The next subsection examines how practitioners made sense of these changes.

%% ================================================================
%% 6.5 sRQ3: Sense-Making of Impacts and Trade-Offs
%% ================================================================
\subsection{sRQ3: Sense-Making of Impacts and Trade-Offs}\label{subsec:srq3}

sRQ3 asks: \textit{``How do EA practitioners make sense of impacts and
trade-offs when adopting GenAI for RA development?''} Verification heuristics
evolved across all three rounds, from personal strategies (Round~1) through
candidate items (Round~2) to refined and rated heuristics (Round~3). Trust
calibration stories and teachability findings follow.

\subsubsection{Verification Heuristics (Rounds~1--2--3)}\label{subsubsec:heuristics}

Verification heuristics evolved across all three rounds, progressively moving
from personal strategies to collectively validated guidance. In Round~1, all
eight participants described individual verification practices, but these varied
widely in formality and specificity. The coded analysis identified five recurring
heuristic patterns: never trust AI-generated numbers, organizational context is
always a human job, treat outputs as hypotheses, verify line-by-line in
compliance contexts, and AI usefulness decreases with domain specificity.

Round~2 formalized these patterns into five candidate heuristics and collected
agreement ratings ($N = 10$). ``Never trust AI-generated numbers'' received the
highest agreement ($M = 4.70$, $SD = 0.67$), followed by ``In compliance
contexts, verify every detail line-by-line'' ($M = 4.20$, $SD = 0.63$). These
two heuristics exhibited both high importance and low variability, indicating
strong consensus.

Round~3 expanded and refined the heuristic set from five to eight items,
incorporating participant-generated additions from the flipboard discussion.
Table~\ref{tab:heuristic-importance} presents the Round~3 importance ratings.
Cross-validation of content (including source verification) received the
highest rating with the lowest variability and no abstentions. Incremental trust
building and context management also rated highly, both above 4.0. Output
structure check received the lowest rating, with one abstention.

\begin{table}[!ht]
    \centering
    \caption{Heuristic Importance Ratings (Round~3, $N = 5$--$7$)}\label{tab:heuristic-importance}
    \begin{tabular}{|p{6.5cm}|c|c|c|c|}
        \hline
        \textbf{Heuristic} & \textbf{$N$} & \textbf{$M$} & \textbf{$SD$} & \textbf{Abst.} \\
        \hline
        Cross-validate content (source verification) & 7 & 4.57 & 0.53 & 0 \\
        \hline
        Incremental trust building & 7 & 4.29 & 1.11 & 0 \\
        \hline
        Context management & 7 & 4.14 & 0.90 & 0 \\
        \hline
        Domain expertise threshold & 7 & 3.86 & 1.07 & 0 \\
        \hline
        Learn through failure & 5 & 3.40 & 0.89 & 2 \\
        \hline
        Contradiction seeding & 6 & 3.33 & 0.82 & 1 \\
        \hline
        Simulate stakeholder reaction & 5 & 3.00 & 1.58 & 2 \\
        \hline
        Output structure check & 6 & 2.83 & 1.33 & 1 \\
        \hline
    \end{tabular}
\end{table}

A tier structure emerged from the ratings. The top tier ($M \geq 4.0$) comprised
cross-validation, incremental trust building, and context management, all
process-oriented heuristics that practitioners can systematically apply. The
middle tier ($M = 3.0$--$3.9$) included domain expertise threshold, learn through
failure, and contradiction seeding, heuristics that depend more on individual
experience. The bottom tier ($M < 3.0$) contained output structure check, rated
lowest and generating the most abstentions. This tier structure showed that process-oriented verification strategies
achieved stronger consensus than experience-dependent ones.

\subsubsection{Trust Calibration Stories (Round~3)}\label{subsubsec:trust-calibration}

Round~3 collected 14 brainstorm responses from seven participants to the trust
calibration prompt (participants could submit multiple responses):
``Describe a specific moment when you had to calibrate trust in AI output.'' Two
themes dominated the responses.

First, domain expertise as prerequisite appeared in multiple contributions. One
participant described validating a solution design where GenAI ``pointed to a
folder that simply didn't exist. One key takeaway here: always verify the
physical components where possible'' (R3 Brainstorm). Another reported that
``a lot of the hallucination issue can only be identified if you know things are
wrong from experience already'' (R3 Brainstorm). A third described an
event-driven architecture where ``the AI output was hallucinating capabilities
of services that did not exist and wrongly explained retry logic for specific
combinations'' (R3 Brainstorm).

Second, the ``junior colleague'' framing recurred. One participant stated:
``Think of AI as a junior architect and yourself as the senior architect'' (R3
Brainstorm). Another described the trust-building process: ``This is very much
teachable: as normally we do this pre-AI when we train or have a buddy system.
Learnt that using the right prompting and almost sharing the elements as if it
were human betters the end output'' (R3 Brainstorm). A third noted that
dependency graph calculations ``were actually better than our manual diagrams.
We missed some links the LLM did not'' (R3 Brainstorm). This indicated that trust calibration was bidirectional; practitioners sometimes
discovered tasks where GenAI outperformed manual approaches, recalibrating trust
upward rather than only guarding against errors.

\subsubsection{Teachability of Calibration (Round~3)}\label{subsubsec:teachability}

The Round~3 flipboard discussion addressed a central tension: if verification
depends on deep domain expertise and personal calibration, can it be taught? The
panel's consensus was nuanced. Calibration is teachable, but it requires domain
expertise as a foundation. One participant explained:

\begin{quote}
Based on experience from the pre-AI era, you had to build a deep understanding
before applying it. If the output of the AI is in line with the principles
learned (scientifically grounded) and experience, at that point I trust it. It
is not the LLM that is in charge. (R3 Brainstorm)
\end{quote}

Several participants identified mechanisms for accelerating calibration in less
experienced practitioners. Cross-checking with authoritative documentation,
having beginners ask the AI to verify its own suggestions, and progressive
exposure from low-stakes to high-stakes tasks all appeared as transferable
strategies. One participant observed that ``the person learning will have another
learning trajectory compared with the persons from the pre-AI era. But in the
end, the knowledge has to be built'' (R3 Brainstorm).

One participant identified a gap in current reinforcement mechanisms. The
observation pointed to the absence of structured feedback loops where
practitioners learn which calibration decisions were correct after the fact.
Without such reinforcement, calibration remains largely experiential, a
limitation relevant to the guidance artifact design.


%% ================================================================
%% 6.6 sRQ4: Systematization into Guidance
%% ================================================================
\subsection{sRQ4: Systematization into Guidance}\label{subsec:srq4}

sRQ4 asks: \textit{``How should effective GenAI practices for RA development be
systematized into guidance for EA practitioners?''} Round~3 rated six candidate
guidance elements, critiqued the draft artifact, and refined the artifact's
structure.

\subsubsection{Guidance Element Priorities (Round~3)}\label{subsubsec:guidance-priorities}

Round~3 presented the draft guidance artifact (v2.5) structured by practice type
and asked participants to rate six candidate elements on importance.
Table~\ref{tab:guidance-priorities} presents the results. Practice description
received the highest rating with the lowest variability, indicating
near-unanimous agreement that describing what each practice is and when it
applies is essential. Verification patterns (practice-specific trust
calibration strategies) received the second-highest rating.

\begin{table}[!ht]
    \centering
    \caption{Guidance Element Priority Ratings (Round~3, $N = 7$)}\label{tab:guidance-priorities}
    \begin{tabular}{|p{6.5cm}|c|c|c|}
        \hline
        \textbf{Guidance Element} & \textbf{$M$} & \textbf{$SD$} & \textbf{Var.\%} \\
        \hline
        Practice description & 4.71 & 0.49 & 22 \\
        \hline
        Verification patterns & 4.14 & 0.69 & 31 \\
        \hline
        Guidance level (directive vs. adaptive) & 3.71 & 0.95 & 44 \\
        \hline
        Experience indicators & 3.43 & 0.98 & 45 \\
        \hline
        Cross-role considerations & 3.43 & 1.27 & 58 \\
        \hline
        Lifecycle phase mapping & 2.43 & 0.53 & 24 \\
        \hline
    \end{tabular}
\end{table}

Lifecycle phase mapping received the lowest priority with low variability,
indicating consensus that mapping practices to specific
lifecycle phases was the least essential element. This finding was noteworthy
given that the conceptual model in Section~\ref{sec:conceptual-model}
emphasizes the RA lifecycle as a structuring lens. The panel's response
suggested that practitioners valued practice descriptions and verification
strategies over lifecycle positioning when consuming guidance.

Cross-role considerations exhibited the highest variability (58\%) among the six
elements, consistent with the contested cross-role dynamics observed in
Section~\ref{subsubsec:cross-role}. This variability reflected genuine
disagreement about whether cross-role considerations belong in practice-level
guidance or at a higher organizational level.

Two additional elements surfaced during the verbal discussion but were not voted
on: criticality-based applicability (applying practices differently based on
system criticality) and proven prompt examples (high-quality prompt templates per
practice). Both were noted for inclusion at the guidance level rather than per
practice.

The SP2 maturity progression results
(Table~\ref{tab:maturity-progression}) further informed the guidance design. The
panel's strong agreement on natural progression ($M = 4.71$) combined with
moderate support for per-category measurement ($M = 3.86$) suggested that
experience indicators should be differentiated per practice category rather than
applied uniformly across the artifact.

\subsubsection{Artifact Critique (Round~3)}\label{subsubsec:artifact-critique}

Round~3 collected ten brainstorm responses from seven participants (participants
could submit multiple responses) critiquing the draft artifact across
four dimensions: structure, format, missing elements, and the
directive-versus-adaptive distinction.

On structure, participants validated organizing by practice type rather than by
lifecycle phase or experience level. One participant stated: ``By practice seems
fine, because if you do by lifecycle, then you will repeat yourself'' (R3
Brainstorm). Another affirmed: ``I think this is the right approach. The decision
on how to apply AI will be done by use case as structured here'' (R3 Brainstorm).

On format, participants expressed interest in multiple consumption modes. Several
mentioned using it as a reference document paired with a checklist. One
participant noted that ``since the guide is already in md format it can be easily
fed to an agent/assistant, so in addition to just a document there can be an
assistant whom you could ask questions about this guide'' (R3 Brainstorm). This
response pointed toward agent-consumable guidance as a format consideration.

On missing elements, participants identified proven prompt examples and additional
context for first-time readers as gaps. One participant noted: ``It would need to
be more verbose. At the moment it assumes a lot of background knowledge'' (R3
Brainstorm). Another suggested including ``proven, high-quality prompt samples'' (R3 Brainstorm)
per practice for less experienced practitioners.

On the directive-versus-adaptive distinction, one participant reported that ``it
took me a few minutes to figure out whether directive and adaptive were in the
text'' (R3 Brainstorm), suggesting the labeling needed clearer presentation.

\subsubsection{Artifact Evolution Across Rounds}\label{subsubsec:artifact-evolution}

The guidance artifact progressed through four stages across the study, following
the emergent artifact approach described in Section~\ref{sec:methodology}.

\textbf{Initial taxonomy} (between Rounds~1--2) emerged from the researcher's
coded analysis of Round~1 data. It organized 19 discovered practices into an
initial taxonomy mapped to RA lifecycle phases.

\textbf{Validated patterns} (between Rounds~2--3) incorporated Round~2 validation
results. The seven practice categories replaced the initial taxonomy. Importance
ratings and variability profiles informed a preliminary directive-versus-adaptive
classification per category.

\textbf{Structured artifact} (presented in Round~3) structured each practice with
six elements: description, lifecycle phases, guidance level, experience
indicators, verification patterns, and cross-role considerations. This
represented the researcher's synthesis of Rounds~1--2 findings into a concrete,
reviewable artifact.

\textbf{Final guidance} (Section~\ref{sec:artifact}) incorporated Round~3
priorities, critique feedback, and the heuristic refinement data into the final
guidance artifact. The panel's priority ratings
(Table~\ref{tab:guidance-priorities}) directly informed which elements received
emphasis.

Four data collection limitations warrant disclosure before the section summary.

%% ================================================================
%% 6.7 Study Limitations Specific to Data Collection
%% ================================================================
\subsection{Study Limitations Specific to Data Collection}\label{subsec:data-limitations}

Four data collection limitations affected the completeness of the findings
presented in this section.

First, Round~3 Step~15 (Future Perspectives brainstorm) was not executed due to
time constraints, and no asynchronous follow-up was conducted. This step was
planned to collect forward-looking data on how GenAI practices in RA development
may evolve. Its absence means the temporal dimension of the conceptual framework
(Section~\ref{subsec:frameworkdimensions}) remains empirically underexplored.

Second, Step~16 (SP2 Maturity Progression vote) was collected asynchronously
after the session rather than during it. All seven Round~3 participants
responded, and the voting data is complete
(Table~\ref{tab:maturity-progression}). Three methodological caveats apply:
(a)~the vote was collected without the preceding group discussion (Step~15) that
the protocol designed as reflective priming, (b)~asynchronous collection
eliminates the real-time deliberation and peer influence that characterize
synchronous Delphi votes, and (c)~participants voted individually rather than
after collectively exploring future perspectives.

Third, Round~3 Steps 13--14 (artifact discussion and guidance element priority
voting) were compressed due to time pressure. While all seven participants voted,
the discussion preceding the vote was shorter than planned. This may have
affected the depth of deliberation informing the priority ratings, though the
voting data itself remains complete.

Fourth, participant counts varied across rounds and across voting items within
rounds. Round~1 had eight participants, Round~2 had ten, and Round~3 had seven.
Within Round~3, abstentions ranged from zero to
two per item, particularly for the heuristic refinement vote. These variations limit the interpretive weight of individual items, though all
items allow meaningful comparison as discovery indicators.

